% TODO checklist: Inspected files:
% - db/postgres_control/README_postgres_control.md
% - db/postgres_control/models/ (all ORM model files)
% - db/postgres_control/alembic/versions/ (migration files 001-026)
% - db/memgraph_domain/README_memgraph_domain.md
% - db/migrations/README_migrations.md
% - Q&A.md (Q10-12: Database layers)

\chapter{Database Schema}
\label{app:database-schema}

This appendix documents the database schemas for the three persistence layers of the Cineca Agentic Platform: PostgreSQL (control plane), Redis (data plane), and Memgraph (graph domain).

%------------------------------------------------------------------------------
\section{PostgreSQL Schema}
\label{app:db:postgresql}

PostgreSQL serves as the authoritative control plane, storing tenants, agents, jobs, providers, models, tools, and audit logs. The schema is managed through 26+ Alembic migrations.

\subsection{Core Tables}

\subsubsection{Tenants Table}

Multi-tenant organization management.

\begin{lstlisting}[language=SQL]
CREATE TABLE tenants (
    id VARCHAR(255) PRIMARY KEY,        -- "tenant-abc123"
    name VARCHAR(255) NOT NULL,         -- "ACME Corporation"
    admin_email VARCHAR(255) NOT NULL,  -- "admin@acme.com" (RFC 5321 max length: 320)
    metadata JSONB NOT NULL DEFAULT '{}', -- {"region": "us-east-1"}
    created_at TIMESTAMPTZ NOT NULL DEFAULT NOW(),
    updated_at TIMESTAMPTZ NOT NULL DEFAULT NOW(),
    version INTEGER NOT NULL DEFAULT 1  -- Optimistic locking
);

-- Constraints
CHECK (char_length(name) BETWEEN 1 AND 255)  -- Name length constraint
UNIQUE INDEX ix_tenants_name_lower_unique ON tenants(LOWER(name));  -- Case-insensitive unique name
INDEX ix_tenants_admin_email_lower ON tenants(LOWER(admin_email));  -- Email lookup index
INDEX ix_tenants_created_at_desc ON tenants(created_at DESC);       -- Pagination index
\end{lstlisting}

\textbf{Constraints:}
\begin{itemize}
    \item \texttt{PRIMARY KEY}: \texttt{id} (VARCHAR, tenant identifier)
    \item \texttt{NOT NULL}: \texttt{id}, \texttt{name}, \texttt{admin_email}, \texttt{metadata}, \texttt{created_at}, \texttt{updated_at}, \texttt{version}
    \item \texttt{UNIQUE}: \texttt{LOWER(name)} (case-insensitive tenant name uniqueness)
    \item \texttt{CHECK}: Name length between 1 and 255 characters
    \item \textbf{Tenant scoping}: Tenant ID serves as isolation boundary; all related tables reference \texttt{tenants.id} via foreign keys
\end{itemize}

\subsubsection{Jobs Table}

Background task management with state machine.

\begin{lstlisting}[language=SQL]
CREATE TABLE jobs (
    id UUID PRIMARY KEY DEFAULT gen_random_uuid(),  -- Requires pgcrypto extension
    type VARCHAR(100) NOT NULL,         -- "demo", "agent.run"
    status VARCHAR(50) NOT NULL,        -- "queued"|"running"|"finished"|"failed"|"cancelled"
    owner_sub VARCHAR(255) NOT NULL,    -- User subject
    tenant_id VARCHAR(255) REFERENCES tenants(id) ON DELETE SET NULL,
    payload_json JSONB,                 -- Input parameters
    result_json JSONB,                  -- Output data
    error_json JSONB,                   -- Error details
    idempotency_key VARCHAR(255),       -- Duplicate prevention
    priority INTEGER NOT NULL DEFAULT 0,
    created_at TIMESTAMPTZ NOT NULL DEFAULT NOW(),
    started_at TIMESTAMPTZ,
    completed_at TIMESTAMPTZ,
    queue_latency_ms INTEGER,
    exec_latency_ms INTEGER,
    etag VARCHAR(64)
);

-- Constraints
UNIQUE INDEX idx_jobs_idempotency ON jobs(owner_sub, idempotency_key) 
    WHERE idempotency_key IS NOT NULL;  -- Partial unique index for idempotency
INDEX idx_jobs_status_created ON jobs(status, created_at DESC);  -- Status query index
INDEX idx_jobs_tenant_created ON jobs(tenant_id, created_at DESC) WHERE tenant_id IS NOT NULL;
\end{lstlisting}

\textbf{Constraints:}
\begin{itemize}
    \item \texttt{PRIMARY KEY}: \texttt{id} (UUID, auto-generated)
    \item \texttt{NOT NULL}: \texttt{id}, \texttt{type}, \texttt{status}, \texttt{owner_sub}, \texttt{priority}, \texttt{created_at}
    \item \texttt{FOREIGN KEY}: \texttt{tenant_id} REFERENCES \texttt{tenants(id)} ON DELETE SET NULL
    \item \texttt{UNIQUE}: \texttt{(owner_sub, idempotency_key)} WHERE \texttt{idempotency_key IS NOT NULL} (partial index for idempotency enforcement)
    \item \textbf{Tenant scoping}: Jobs reference \texttt{tenant_id} for multi-tenant isolation; nullable to support global/system jobs
    \item \textbf{Note}: \texttt{gen_random_uuid()} requires \texttt{pgcrypto} extension
\end{itemize}

\subsection{Agent Tables}

\subsubsection{Agent Sessions}

Stateful conversation context.

\begin{lstlisting}[language=SQL]
CREATE TABLE agent_sessions (
    session_id UUID PRIMARY KEY DEFAULT gen_random_uuid(),
    user_id VARCHAR(255) NOT NULL,
    tenant_id VARCHAR(255) REFERENCES tenants(id) ON DELETE SET NULL,
    status VARCHAR(50) NOT NULL,        -- "active"|"completed"|"cancelled"|"failed"
    manager VARCHAR(100),               -- Manager type
    temperature FLOAT NOT NULL DEFAULT 0.2,
    max_steps INTEGER NOT NULL DEFAULT 8,
    tools JSONB NOT NULL DEFAULT '[]',  -- Enabled tools
    session_metadata JSONB NOT NULL DEFAULT '{}',
    last_step_id UUID,
    last_step_seq INTEGER NOT NULL DEFAULT 0,
    etag VARCHAR(64),
    created_at TIMESTAMPTZ NOT NULL DEFAULT NOW(),
    updated_at TIMESTAMPTZ NOT NULL DEFAULT NOW()
);

-- Constraints
INDEX idx_sessions_user_tenant ON agent_sessions(user_id, tenant_id);  -- User/tenant lookup
\end{lstlisting}

\textbf{Constraints:}
\begin{itemize}
    \item \texttt{PRIMARY KEY}: \texttt{session_id} (UUID)
    \item \texttt{NOT NULL}: \texttt{session_id}, \texttt{user_id}, \texttt{status}, \texttt{temperature}, \texttt{max_steps}, \texttt{tools}, \texttt{session_metadata}, \texttt{last_step_seq}, \texttt{created_at}, \texttt{updated_at}
    \item \texttt{FOREIGN KEY}: \texttt{tenant_id} REFERENCES \texttt{tenants(id)} ON DELETE SET NULL
    \item \textbf{Tenant scoping}: Sessions are tenant-scoped via \texttt{tenant_id} foreign key
\end{itemize}

\subsubsection{Agent Steps}

Individual steps within sessions.

\begin{lstlisting}[language=SQL]
CREATE TABLE agent_steps (
    step_id UUID PRIMARY KEY DEFAULT gen_random_uuid(),
    session_id UUID NOT NULL REFERENCES agent_sessions(session_id) ON DELETE CASCADE,
    seq INTEGER NOT NULL,               -- Sequence number
    type VARCHAR(50) NOT NULL,          -- "user"|"assistant"|"tool"|"system"|"error"|"message"
    message TEXT,
    tool VARCHAR(255),
    input JSONB,
    output JSONB,
    status VARCHAR(50),                 -- "queued"|"running"|"completed"|"failed"|"cancelled"
    started_at TIMESTAMPTZ,
    finished_at TIMESTAMPTZ,
    created_at TIMESTAMPTZ NOT NULL DEFAULT NOW()
);

-- Constraints
INDEX idx_steps_session_seq ON agent_steps(session_id, seq);  -- Composite index for session step ordering
INDEX idx_steps_timestamps ON agent_steps(started_at, finished_at) 
    WHERE started_at IS NOT NULL;  -- Partial index for active/running steps
UNIQUE INDEX idx_steps_session_seq_unique ON agent_steps(session_id, seq);  -- Ensure sequential ordering per session
\end{lstlisting}

\textbf{Constraints:}
\begin{itemize}
    \item \texttt{PRIMARY KEY}: \texttt{step_id} (UUID)
    \item \texttt{NOT NULL}: \texttt{step_id}, \texttt{session_id}, \texttt{seq}, \texttt{type}, \texttt{created_at}
    \item \texttt{FOREIGN KEY}: \texttt{session_id} REFERENCES \texttt{agent_sessions(session_id)} ON DELETE CASCADE (cascading delete)
    \item \texttt{UNIQUE}: \texttt{(session_id, seq)} (ensures sequential ordering within each session)
    \item \textbf{Tenant scoping}: Steps inherit tenant scope from parent session via \texttt{session_id} foreign key
\end{itemize}

\subsubsection{Agent Runs}

Single agent execution records.

\begin{lstlisting}[language=SQL]
CREATE TABLE agent_runs (
    run_id UUID PRIMARY KEY DEFAULT gen_random_uuid(),
    session_id UUID REFERENCES agent_sessions(session_id),
    user_id VARCHAR(255) NOT NULL,
    tenant_id VARCHAR(255) REFERENCES tenants(id),
    model_instance_name VARCHAR(255),
    model_id VARCHAR(255),
    provider_name VARCHAR(255),
    provider_id VARCHAR(255),
    status VARCHAR(50) NOT NULL,        -- "queued"|"running"|"succeeded"|"failed"|"cancelled"
    latency_ms INTEGER,
    trace_id VARCHAR(255),
    request_id VARCHAR(255),
    llm_error_type VARCHAR(100),
    llm_error_message TEXT,
    llm_error_occurred_at TIMESTAMPTZ,
    todos JSONB DEFAULT '[]',
    steps JSONB DEFAULT '[]',
    output JSONB,
    warnings JSONB DEFAULT '[]',
    metrics JSONB,
    run_metadata JSONB,
    created_at TIMESTAMPTZ DEFAULT NOW(),
    started_at TIMESTAMPTZ,
    finished_at TIMESTAMPTZ
);

CREATE INDEX idx_runs_tenant_user ON agent_runs(tenant_id, user_id, started_at);
CREATE INDEX idx_runs_status ON agent_runs(status, started_at);
CREATE INDEX idx_runs_output_gin ON agent_runs USING gin(output);
CREATE INDEX idx_runs_metrics_gin ON agent_runs USING gin(metrics);
\end{lstlisting}

\subsection{Provider Tables}

\subsubsection{Providers}

LLM provider configuration.

\begin{lstlisting}[language=SQL]
CREATE TABLE providers (
    id VARCHAR(255) PRIMARY KEY,        -- "ollama-local"
    name VARCHAR(255) NOT NULL,
    type VARCHAR(100) NOT NULL,         -- "openai_compatible", "ollama"
    base_url VARCHAR(1000),
    model VARCHAR(255),
    tenant_id VARCHAR(255) REFERENCES tenants(id) ON DELETE CASCADE,  -- NULL for global
    config_json JSONB NOT NULL DEFAULT '{}',
    has_api_key BOOLEAN NOT NULL DEFAULT FALSE,
    enabled BOOLEAN NOT NULL DEFAULT TRUE,
    created_at TIMESTAMPTZ NOT NULL DEFAULT NOW(),
    updated_at TIMESTAMPTZ NOT NULL DEFAULT NOW()
);

-- Constraints
UNIQUE INDEX uq_provider_tenant_name ON providers(tenant_id, name);  -- Tenant-scoped unique provider name
\end{lstlisting}

\textbf{Constraints:}
\begin{itemize}
    \item \texttt{PRIMARY KEY}: \texttt{id} (VARCHAR)
    \item \texttt{NOT NULL}: \texttt{id}, \texttt{name}, \texttt{type}, \texttt{config_json}, \texttt{has_api_key}, \texttt{enabled}, \texttt{created_at}, \texttt{updated_at}
    \item \texttt{FOREIGN KEY}: \texttt{tenant_id} REFERENCES \texttt{tenants(id)} ON DELETE CASCADE (nullable for global providers)
    \item \texttt{UNIQUE}: \texttt{(tenant_id, name)} (provider name uniqueness per tenant; global providers have \texttt{tenant_id=NULL})
    \item \textbf{Tenant scoping}: Provider uniqueness is enforced per tenant via unique index; \texttt{tenant_id=NULL} represents global providers shared across all tenants
\end{itemize}

\subsubsection{Provider Secrets}

Encrypted API keys (separate table for security).

\begin{lstlisting}[language=SQL]
CREATE TABLE provider_secrets (
    provider_id VARCHAR(255) PRIMARY KEY REFERENCES providers(id) ON DELETE CASCADE,
    api_key_encrypted TEXT NOT NULL,    -- Fernet-encrypted
    created_at TIMESTAMPTZ NOT NULL DEFAULT NOW(),
    updated_at TIMESTAMPTZ NOT NULL DEFAULT NOW()
);
\end{lstlisting}

\textbf{Constraints:}
\begin{itemize}
    \item \texttt{PRIMARY KEY}: \texttt{provider_id} (VARCHAR, also foreign key)
    \item \texttt{NOT NULL}: \texttt{provider_id}, \texttt{api_key_encrypted}, \texttt{created_at}, \texttt{updated_at}
    \item \texttt{FOREIGN KEY}: \texttt{provider_id} REFERENCES \texttt{providers(id)} ON DELETE CASCADE (one-to-one relationship)
    \item \textbf{Security}: Secrets stored in separate table with encrypted values; deletion cascades when provider is deleted
\end{itemize}

\subsection{Model Tables}

\subsubsection{Model Instances}

Loaded LLM model instances.

\begin{lstlisting}[language=SQL]
CREATE TABLE model_instances (
    id UUID PRIMARY KEY DEFAULT gen_random_uuid(),
    instance_name VARCHAR(255) NOT NULL,
    provider_id VARCHAR(255) REFERENCES providers(id) ON DELETE SET NULL,
    model_id VARCHAR(255) NOT NULL,     -- Provider-specific model ID
    tenant_id VARCHAR(255) REFERENCES tenants(id) ON DELETE SET NULL,
    enabled BOOLEAN NOT NULL DEFAULT TRUE,
    loaded BOOLEAN NOT NULL DEFAULT FALSE,
    capabilities JSONB NOT NULL DEFAULT '["chat"]',
    config_json JSONB NOT NULL DEFAULT '{}',
    created_at TIMESTAMPTZ NOT NULL DEFAULT NOW(),
    updated_at TIMESTAMPTZ NOT NULL DEFAULT NOW()
);

-- Constraints
INDEX idx_instances_provider ON model_instances(provider_id);
INDEX idx_instances_tenant ON model_instances(tenant_id) WHERE tenant_id IS NOT NULL;
UNIQUE INDEX idx_instances_tenant_name ON model_instances(tenant_id, instance_name) 
    WHERE tenant_id IS NOT NULL;  -- Tenant-scoped unique instance name
\end{lstlisting}

\textbf{Constraints:}
\begin{itemize}
    \item \texttt{PRIMARY KEY}: \texttt{id} (UUID)
    \item \texttt{NOT NULL}: \texttt{id}, \texttt{instance_name}, \texttt{model_id}, \texttt{enabled}, \texttt{loaded}, \texttt{capabilities}, \texttt{config_json}, \texttt{created_at}, \texttt{updated_at}
    \item \texttt{FOREIGN KEY}: \texttt{provider_id} REFERENCES \texttt{providers(id)} ON DELETE SET NULL; \texttt{tenant_id} REFERENCES \texttt{tenants(id)} ON DELETE SET NULL
    \item \texttt{UNIQUE}: \texttt{(tenant_id, instance_name)} WHERE \texttt{tenant_id IS NOT NULL} (partial unique index for tenant-scoped instance names)
    \item \textbf{Tenant scoping}: Instance names must be unique per tenant; global instances have \texttt{tenant_id=NULL}
\end{itemize}

\subsubsection{User Default Models}

Per-user model preferences.

\begin{lstlisting}[language=SQL]
CREATE TABLE user_default_models (
    id UUID PRIMARY KEY DEFAULT gen_random_uuid(),  -- Changed from SERIAL to UUID
    user_id VARCHAR(255) NOT NULL,
    tenant_id VARCHAR(255) NOT NULL,
    chat_instance_id UUID NOT NULL REFERENCES model_instances(id) ON DELETE CASCADE,
    created_by VARCHAR(255) NOT NULL,
    created_at TIMESTAMPTZ NOT NULL DEFAULT NOW(),
    updated_at TIMESTAMPTZ NOT NULL DEFAULT NOW(),
    UNIQUE(user_id, tenant_id)  -- One default model per user per tenant
);

-- Constraints
INDEX idx_user_default_models_user_tenant ON user_default_models(user_id, tenant_id);
INDEX idx_user_default_models_instance ON user_default_models(chat_instance_id);
\end{lstlisting}

\textbf{Constraints:}
\begin{itemize}
    \item \texttt{PRIMARY KEY}: \texttt{id} (UUID)
    \item \texttt{NOT NULL}: \texttt{id}, \texttt{user_id}, \texttt{tenant_id}, \texttt{chat_instance_id}, \texttt{created_by}, \texttt{created_at}, \texttt{updated_at}
    \item \texttt{FOREIGN KEY}: \texttt{chat_instance_id} REFERENCES \texttt{model_instances(id)} ON DELETE CASCADE
    \item \texttt{UNIQUE}: \texttt{(user_id, tenant_id)} (one default model preference per user per tenant)
    \item \textbf{Tenant scoping}: User defaults are tenant-scoped via \texttt{tenant_id}; uniqueness enforced per (user, tenant) pair
\end{itemize}

\subsection{Tool Tables}

\begin{lstlisting}[language=SQL]
CREATE TABLE tools (
    id VARCHAR(255) PRIMARY KEY,
    name VARCHAR(255) NOT NULL,
    version VARCHAR(50) NOT NULL,
    description TEXT,
    input_schema JSONB,
    output_schema JSONB,
    owner_tenant_id VARCHAR(255) REFERENCES tenants(id) ON DELETE SET NULL,
    version_number INTEGER NOT NULL DEFAULT 1,
    created_at TIMESTAMPTZ NOT NULL DEFAULT NOW(),
    updated_at TIMESTAMPTZ NOT NULL DEFAULT NOW()
);

CREATE TABLE tool_invocations (
    eid UUID PRIMARY KEY DEFAULT gen_random_uuid(),
    tool_name VARCHAR(255) NOT NULL,
    status VARCHAR(50),
    params_json JSONB,
    result_json JSONB,
    user_id VARCHAR(255),
    tenant_id VARCHAR(255) REFERENCES tenants(id) ON DELETE SET NULL,
    latency_ms INTEGER,
    created_at TIMESTAMPTZ NOT NULL DEFAULT NOW()
);

-- Constraints
INDEX idx_tools_owner_tenant ON tools(owner_tenant_id) WHERE owner_tenant_id IS NOT NULL;
INDEX idx_tool_invocations_tenant_user ON tool_invocations(tenant_id, user_id, created_at DESC);
INDEX idx_tool_invocations_tool_tenant ON tool_invocations(tool_name, tenant_id, created_at DESC);
\end{lstlisting}

\textbf{Constraints:}
\begin{itemize}
    \item \textbf{Tools table}:
    \begin{itemize}
        \item \texttt{PRIMARY KEY}: \texttt{id} (VARCHAR)
        \item \texttt{NOT NULL}: \texttt{id}, \texttt{name}, \texttt{version}, \texttt{version_number}, \texttt{created_at}, \texttt{updated_at}
        \item \texttt{FOREIGN KEY}: \texttt{owner_tenant_id} REFERENCES \texttt{tenants(id)} ON DELETE SET NULL (nullable for global tools)
        \item \textbf{Tenant scoping}: Tools can be tenant-owned or global (\texttt{owner_tenant_id=NULL})
    \end{itemize}
    \item \textbf{Tool Invocations table}:
    \begin{itemize}
        \item \texttt{PRIMARY KEY}: \texttt{eid} (UUID)
        \item \texttt{NOT NULL}: \texttt{eid}, \texttt{tool_name}, \texttt{created_at}
        \item \texttt{FOREIGN KEY}: \texttt{tenant_id} REFERENCES \texttt{tenants(id)} ON DELETE SET NULL
        \item \textbf{Tenant scoping}: Invocations are tenant-scoped for audit and isolation
    \end{itemize}
\end{itemize}

\subsection{Audit Tables}

\begin{lstlisting}[language=SQL]
CREATE TABLE internal_ops_events (
    id BIGSERIAL PRIMARY KEY,
    correlation_id VARCHAR(255) NOT NULL,
    event_type VARCHAR(100) NOT NULL,
    actor_sub VARCHAR(255) NOT NULL,
    actor_type VARCHAR(50),
    endpoint VARCHAR(255) NOT NULL,
    http_method VARCHAR(10) NOT NULL,
    request_params JSONB,
    request_body JSONB,
    http_status INTEGER NOT NULL,
    response_body JSONB,
    operation_result VARCHAR(50) NOT NULL,
    duration_ms INTEGER,
    idempotency_key VARCHAR(255),
    is_idempotency_replay BOOLEAN NOT NULL DEFAULT FALSE,
    cache_status VARCHAR(20),
    error_type VARCHAR(100),
    error_message TEXT,
    created_at TIMESTAMPTZ NOT NULL DEFAULT NOW(),
    metadata JSONB
);

-- Constraints
INDEX idx_ops_events_correlation ON internal_ops_events(correlation_id);  -- Trace lookup
INDEX idx_ops_events_actor_time ON internal_ops_events(actor_sub, created_at DESC);  -- Actor audit trail
INDEX idx_ops_events_created_at_desc ON internal_ops_events(created_at DESC);  -- Time-based queries
\end{lstlisting}

\textbf{Constraints:}
\begin{itemize}
    \item \texttt{PRIMARY KEY}: \texttt{id} (BIGSERIAL, auto-incrementing)
    \item \texttt{NOT NULL}: \texttt{id}, \texttt{correlation_id}, \texttt{event_type}, \texttt{actor_sub}, \texttt{endpoint}, \texttt{http_method}, \texttt{http_status}, \texttt{operation_result}, \texttt{is_idempotency_replay}, \texttt{created_at}
    \item \textbf{Tenant scoping}: Audit events capture tenant context via \texttt{request_params} or \texttt{metadata} JSONB fields (application-level enforcement)
\end{itemize}

%------------------------------------------------------------------------------
\section{Memgraph Schema}
\label{app:db:memgraph}

Memgraph stores the bioinformatics knowledge graph with users, institutions, tasks, and files.

\subsection{Node Types}

\begin{table}[ht]
\centering
\caption{Memgraph Node Labels}
\label{tab:memgraph-nodes}
\small
\begin{tabularx}{\textwidth}{lX}
    \hline
    \textbf{Label} & \textbf{Key Properties} \\
    \hline
    \texttt{User} & \texttt{user\_id}, \texttt{firstName}, \texttt{lastName}, \texttt{email}, \texttt{user\_name} \\
    \texttt{Institution} & \texttt{name} \\
    \texttt{Blast} & \texttt{task\_id}, \texttt{blasttype}, \texttt{blast\_version}, \texttt{dbname}, \texttt{status} \\
    \texttt{CreateDb} & \texttt{task\_id}, \texttt{dbtype}, \texttt{dbname}, \texttt{status} \\
    \texttt{SearchbyTaxon} & \texttt{task\_id}, \texttt{taxon}, \texttt{tool}, \texttt{status} \\
    \texttt{File} & \texttt{file\_id}, \texttt{user\_filename}, \texttt{size}, \texttt{extension} \\
    \texttt{Fasta} & \texttt{file\_id}, \texttt{user\_filename}, \texttt{size} \\
    \texttt{BlastDb} & \texttt{file\_id}, \texttt{dbname}, \texttt{size} \\
    \hline
\end{tabularx}
\end{table}

\subsection{Relationship Types}

\begin{table}[ht]
\centering
\caption{Memgraph Relationship Types}
\label{tab:memgraph-rels}
\small
\begin{tabularx}{\textwidth}{llX}
    \hline
    \textbf{Type} & \textbf{Direction} & \textbf{Description} \\
    \hline
    \texttt{WORKS\_AT} & \texttt{(User)$\rightarrow$(Institution)} & User employment \\
    \texttt{RUNS} & \texttt{(User)$\rightarrow$(Task)} & User executes task \\
    \texttt{INPUT} & \texttt{(File)$\rightarrow$(Task)} & File input to task \\
    \texttt{OUTPUT} & \texttt{(Task)$\rightarrow$(File)} & Task produces file \\
    \hline
\end{tabularx}
\end{table}

\subsection{Schema Diagram}

\begin{verbatim}
                    +--------------+
                    | Institution  |
                    +------^-------+
                           | WORKS_AT
                    +------+-------+
                    |    User      |
                    +------+-------+
                           | RUNS
           +---------------+----------------+
           |               |                |
    +------v------+ +------v------+ +-------v-------+
    |    Blast    | | CreateDb    | | SearchbyTaxon |
    +------+------+ +------+------+ +-------+-------+
           |               |                |
    INPUT -+- OUTPUT INPUT -+- OUTPUT INPUT -+- OUTPUT
           |               |                |
    +------v------+ +------v------+ +-------v-------+
    |    File     | |  BlastDb    | |     Fasta     |
    +-------------+ +-------------+ +---------------+
\end{verbatim}

%------------------------------------------------------------------------------
\section{Redis Key Patterns}
\label{app:db:redis}

Redis serves as the low-latency data plane for caching, job queues, rate limiting, and session state.

\begin{table}[ht]
\centering
\caption{Redis Key Patterns}
\label{tab:redis-keys}
\small
\begin{tabularx}{\textwidth}{lrX}
    \hline
    \textbf{Pattern} & \textbf{TTL} & \textbf{Description} \\
    \hline
    \texttt{job:\{id\}:status} & 10d & Job state and metadata \\
    \texttt{job:\{id\}:events} & 10d & Job event ring buffer \\
    \texttt{job:queue:\{type\}} & -- & Job queue (BRPOP) \\
    \texttt{idem:\{key\}} & 24h & Idempotency cache \\
    \texttt{rate:\{user\}:\{window\}} & -- & Rate limit counters (ZSET) \\
    \texttt{session:\{id\}} & 1h & Session state cache \\
    \texttt{provider:\{id\}:health} & 2h & Provider health status \\
    \texttt{default:model:\{scope\}} & 15m & Default model cache \\
    \texttt{jwks:\{issuer\}} & 1h & JWKS public keys \\
    \texttt{circuit:\{provider\}} & -- & Circuit breaker state \\
    \hline
\end{tabularx}
\end{table}

%------------------------------------------------------------------------------
\section{Naming Conventions and Extensions}
\label{app:db:naming}

\subsection{Index Naming}

The schema uses the following index naming conventions:

\begin{itemize}
    \item \texttt{ix\_*}: Standard indexes (e.g., \texttt{ix\_tenants\_name\_lower\_unique})
    \item \texttt{idx\_*}: Performance indexes (e.g., \texttt{idx\_jobs\_status\_created})
    \item \texttt{uq\_*}: Unique indexes (e.g., \texttt{uq\_provider\_tenant\_name})
    \item Unique indexes may also use \texttt{ix\_*} prefix with \texttt{\_unique} suffix
\end{itemize}

\subsection{UUID Generation}

The schema uses PostgreSQL's \texttt{gen\_random\_uuid()} function for UUID primary keys. This requires the \texttt{pgcrypto} extension:

\begin{lstlisting}[language=SQL]
CREATE EXTENSION IF NOT EXISTS pgcrypto;
\end{lstlisting}

Alternatively, the \texttt{uuid-ossp} extension can be used with \texttt{uuid\_generate\_v4()}:

\begin{lstlisting}[language=SQL]
CREATE EXTENSION IF NOT EXISTS "uuid-ossp";
-- Then use: DEFAULT uuid_generate_v4()
\end{lstlisting}

The current implementation uses \texttt{pgcrypto} for UUID generation.

\subsection{Constraint Naming}

\begin{itemize}
    \item \textbf{Primary keys}: Use column name as constraint name (implicit \texttt{PRIMARY KEY})
    \item \textbf{Foreign keys}: Use descriptive names when explicitly named (e.g., \texttt{FK\_jobs\_tenant\_id})
    \item \textbf{Check constraints}: Use \texttt{ck\_*} prefix (e.g., \texttt{ck\_tenants\_name\_length})
    \item \textbf{Unique constraints}: Use \texttt{uq\_*} prefix or \texttt{ix\_*\_unique} for unique indexes
\end{itemize}

%------------------------------------------------------------------------------
\section{Migration History}
\label{app:db:migrations}

\begin{table}[ht]
\centering
\caption{Alembic Migration History}
\label{tab:migrations}
\small
\begin{tabularx}{\textwidth}{lX}
    \hline
    \textbf{Revision} & \textbf{Description} \\
    \hline
    001 & Initial tenants table \\
    002 & Tools and invocations tables \\
    003 & Jobs and job events tables \\
    004 & Providers and secrets tables \\
    005 & Built-in manifests tables \\
    006 & Model instances tables \\
    007 & User default models \\
    008 & Agent tables (sessions, steps, runs) \\
    009--015 & Session and run column additions \\
    016--019 & Model defaults fixes \\
    020--025 & Agent run enhancements \\
    026 & Performance indexes \\
    \hline
\end{tabularx}
\end{table}

